\documentclass[a4paper,12pt]{article}
\usepackage[utf8]{inputenc}
\usepackage[brazilian]{babel}
\usepackage[T1]{fontenc}
\usepackage{geometry}
\geometry{top=2.5cm, bottom=2.5cm, left=3cm, right=3cm}
\usepackage{titlesec}
\usepackage{enumitem}
\usepackage{hyperref}
\usepackage{xcolor}
\usepackage{tcolorbox}

\title{\textbf{Resumo: Resolução CONSUP/IFSC nº 142/2025}}
\author{Comissão de Reformulação dos PPCs}
\date{\today}

\begin{document}

\maketitle

\section{Introdução}
Este documento resume os principais pontos da \textbf{Resolução CONSUP/IFSC nº 142, de 27 de março de 2025}, que estabelece as novas Diretrizes Curriculares para os Cursos Técnicos Integrados ao Ensino Médio do Instituto Federal de Santa Catarina (IFSC). Esta normativa visa fortalecer a integração entre a formação geral e a profissional, garantindo princípios de formação humana integral.

\section{Estrutura Curricular (Art. 7º)}
Os currículos devem ser organizados em três grandes blocos integrados:
\begin{itemize}
    \item \textbf{Formação Geral:} Ciências Humanas, Ciências da Natureza, Linguagens (Português, Inglês, Artes, Educação Física e obrigatoriamente Espanhol ou Libras) e Matemática.
    \item \textbf{Formação Técnica:} Conforme definido no Catálogo Nacional de Cursos Técnicos (CNCT).
    \item \textbf{Núcleo Politécnico Comum:} Espaço de integração transversal (Projetos Integradores, Oficinas de Integração, etc.).
\end{itemize}

\section{Diretrizes de Carga Horária (Pontos Críticos)}
A resolução define limites claros que devem ser observados na reformulação dos PPCs:

\begin{tcolorbox}[colback=blue!5!white,colframe=blue!75!black,title=Cargas Horárias Mínimas]
\begin{description}[leftmargin=!,labelwidth=5cm]
    \item[Formação Geral:] Mínimo de \textbf{2.100 horas} (Art. 8º, I).
    \item[Núcleo Politécnico:] Mínimo de \textbf{120 horas} (Art. 7º, III, b).
    \item[Unidade Curricular:] Mínimo de \textbf{40 horas} semestrais (Art. 12º).
    \item[Componente Curricular (Soma):] Mínimo de \textbf{120 horas} ao longo do curso para componentes da Formação Geral (Art. 9º).
\end{description}
\end{tcolorbox}

\subsection{Carga Horária Total do Curso (Art. 10º)}
A carga horária total é determinada pela carga horária da habilitação técnica (800h, 1.000h ou 1.200h):
\begin{itemize}
    \item Habilitação de 800h $\rightarrow$ Total Mínimo: \textbf{3.000 horas}.
    \item Habilitação de 1.000h $\rightarrow$ Total Mínimo: \textbf{3.100 horas}.
    \item Habilitação de 1.200h $\rightarrow$ Total Mínimo: \textbf{3.300 horas}.
\end{itemize}
\textit{Nota: O limite máximo pode extrapolar em até 10\% esses valores (conforme RDP).}

\section{Outras Disposições Importantes}
\begin{itemize}
    \item \textbf{Duração:} Os cursos devem ter duração entre \textbf{3 e 4 anos} (Art. 14º).
    \item \textbf{Língua Estrangeira:} É obrigatória a oferta de Inglês e, no mínimo, uma opção entre Espanhol ou Libras (Art. 7º, §1º).
    \item \textbf{Gestão Coletiva:} Possibilidade de criação de \textbf{Núcleos Pedagógicos Estruturantes} para a gestão do curso (Art. 41º).
    \item \textbf{Encontro Estadual:} Realização bienal do Encontro Estadual do Ensino Médio Integrado (Art. 42º).
\end{itemize}

\section{Conclusão para a Reunião}
Para a reunião de hoje, os pontos de maior impacto para o curso de \textbf{Informática} e \textbf{Administração} serão a necessidade de ajuste da Formação Geral para o mínimo de 2.100h e a institucionalização das 120h de Projetos Integradores (Núcleo Politécnico).

\end{document}
